% §1.3 A First Look at the Boolean Universe
%
% Cross-reference resolution status:
%
% ALREADY DEFINED (within chapter 1, by this section):
%   \label{sec:boolean-universe-first-look}  -- this section
%
% RESOLVED (already in source tree at Plan 4 Task 4 HEAD):
%   \Cref{sec:three-faces}              -- §1.1 (drafted)
%   \Cref{sec:quantifying-noise}        -- §1.2 (drafted)
%   \Cref{sec:bernoulli-bool-formal}    -- chapter 4 §4.5 (drafted)
%   \Cref{def:bernoulli-bool}           -- chapter 4 §4.5 (drafted)
%
% FORWARD REFS to later chapter 1 sections (resolve in Plan 4 Task 5):
%   \Cref{sec:sampling-rates}           -- §1.4 (Task 5)
%
% This section walks the four-cell 2x2 outcome matrix of a noisy bit
% qualitatively, names the smallest non-trivial universe, and defers the
% joint PMF and channel matrix to §4.5. Defining this label retires the
% two outstanding forward references in ch04/bernoulli-bool-formal.tex
% lines 129 and 267 (and the one in the closing paragraph at line 353).

\section{A First Look at the Boolean Universe}
\label{sec:boolean-universe-first-look}

A noisy bit has two possible inputs and two possible outputs. The truth on
any given trial is either $\top$ or $\bot$, and the reported value is also
either $\top$ or $\bot$. Cross those two binary choices and four
input-output combinations exhaust the joint behavior of the device. We
arrange the combinations in a $2 \times 2$ outcome matrix indexed by the
true input (rows) and the observed report (columns):

\begin{center}
\begin{tabular}{l|cc}
\hline
              & \textbf{report $\top$} & \textbf{report $\bot$} \\
\hline
\textbf{truth $\top$} & correct positive & false negative \\
\textbf{truth $\bot$} & false positive   & correct negative \\
\hline
\end{tabular}
\end{center}

\noindent The four cells are the universe of joint behavior. Anything the
noisy bit can do on a single trial is one of these four outcomes; there
are no others. The chapter from here on is about the structure of these
four cells, what populates each one, how the populations move when the
device is composed with another, and how the rates of \Cref{sec:quantifying-noise} fit on top of them.

\subsection*{The four cells, narratively}

Take the cells one at a time.

\emph{Truth $\top$, report $\top$.} The input was $\top$ and the device
reported $\top$. Call this a \emph{correct positive}: the noisy bit got
it right on a $\top$-input. In the spam-filter vocabulary of
\Cref{sec:three-faces} this is a spam message correctly flagged. In the
biased-coin face this is the in-handheld trial that came up heads as it
usually does. In the binary-channel face this is the bit that traveled
the wire without flipping.

\emph{Truth $\top$, report $\bot$.} The input was $\top$ and the device
reported $\bot$. This is the \emph{false negative}: the noisy bit
missed. The spam filter dropped a spam message into the inbox. The
coin came up tails on a trial that should usually have produced heads.
The wire flipped a transmitted $\top$ into a $\bot$. Whatever the
mechanism, the system replied with the wrong answer in the
$\top$-to-$\bot$ direction.

\emph{Truth $\bot$, report $\top$.} The input was $\bot$ and the device
reported $\top$. This is the \emph{false positive}: the noisy bit
fired when it should have stayed silent. The spam filter flagged a
legitimate message. The coin came up heads on a trial when heads was
unlikely. The wire flipped a $\bot$ into a $\top$.

\emph{Truth $\bot$, report $\bot$.} The input was $\bot$ and the device
reported $\bot$. The \emph{correct negative}: the noisy bit stayed
silent when it should have. The ham message passed through, the
unlikely-heads trial produced tails as expected, the $\bot$-bit
traveled the wire unchanged.

Those are the four cells. Their names are exactly the four labels of
the contingency table in \Cref{sec:quantifying-noise}: TP, FN, FP, TN.
This is not a coincidence; it is the same four-way split seen from a
different angle. The contingency table of \Cref{sec:quantifying-noise}
records what happens when we run the noisy bit $n$ times and \emph{count}
how many trials land in each cell. The outcome matrix above records the
\emph{cells themselves}, abstracted from any particular run. The four
cells of the joint distribution and the four cells of the empirical
contingency table are the same four cells. The empirical table is a
sample drawn from the joint distribution over those cells.

\subsection*{Why this is a ``universe''}

The input domain $\{\top, \bot\}$ has exactly two elements. We will call
this set the \emph{Boolean universe} and write it $U = \{\top, \bot\}$.
The word \emph{universe} is borrowed from set-theoretic usage: it is the
ambient set out of which subsets are formed and over which predicates are
evaluated. For our purposes the Boolean universe is the smallest
non-trivial universe of interest. With one element, no nontrivial
predicate exists. With two, we can already ask whether an input is one
specific value or the other, which is the simplest discrimination a
classifier can perform.

The smallness is the point. Part~I of this book develops the algebra of
this universe and nothing larger. Every result, every composition rule,
every error-rate formula in the chapters that follow is a statement
about predicates whose input is a single Boolean. The reason for that
restriction is pedagogical and structural: the Boolean universe is the
\emph{atom} of the framework, and every richer universe (multi-bit
inputs, finite enumerations, arbitrary algebraic types) is built up
from Booleans by Kronecker products in Part~II and beyond. If we
understand the algebra of the four cells, we understand a great deal
about systems that look more complicated on the surface but reduce, on
inspection, to compositions of noisy bits.

\subsection*{The deferral}

A fully formal description of the four cells would assign each one a
probability mass. Given the rates of \Cref{sec:quantifying-noise} and
the prevalence of $\top$ in the input stream, the four cells receive
four numbers that sum to one. The resulting object is a joint
distribution over the four pairs (truth, report). Multiplying it out
involves the rates $\fprate$, $\fnrate$, $\tprate$, $\tnrate$, and the
prior probability of a $\top$-input. We will not write that joint
distribution here.

The reason for the deferral is structural, not coy. The formal definition
of the object we are circling, the type whose instances are exactly the
noisy bits with two free rate parameters, lives in
\Cref{sec:bernoulli-bool-formal}. That section earns the formal step by
first introducing the two axioms that make the type well-posed (in
chapter~4 \S\S4.2 and 4.3) and then proving the Kronecker factorization
that justifies the parameter economy (chapter~4 \S4.4). Writing the joint
PMF in §1.3 would either repeat that machinery or assert it without
support. Both are bad pedagogy. The honest path is to acknowledge that
the four cells admit a joint distribution, name where that distribution
gets written down, and continue. \Cref{sec:bernoulli-bool-formal}
returns to this picture with the formal definition; the table of joint
probabilities, the channel matrix, and the worked biased-coin instance
all appear there.

\subsection*{A qualitative example: the biased coin}

The biased coin from \Cref{sec:three-faces} runs through the book, and it
is worth seeing what the four cells look like for a specific instance.
Take $\fprate = 0.1$ and $\fnrate = 0.2$: the coin reports a $\top$ when
the truth is $\bot$ one trial in ten on average, and reports a $\bot$
when the truth is $\top$ one trial in five on average.

In the outcome matrix this means the following. Of every five
$\top$-trials, roughly four land in the \emph{correct positive} cell and
roughly one lands in the \emph{false negative} cell: the true-positive
rate is $\tprate = 1 - \fnrate = 0.8$ and the false-negative rate is
$\fnrate = 0.2$. Of every ten $\bot$-trials, nine land in the \emph{correct
negative} cell and one lands in the \emph{false positive} cell:
$\tnrate = 1 - \fprate = 0.9$ and $\fprate = 0.1$.

These four conditional shares (four-fifths, one-fifth, one-tenth,
nine-tenths) say what fraction of each \emph{row} the noisy bit places in
each \emph{column}. They do not say what fraction of all trials lands in
each cell, because that fraction also depends on how many trials had
truth $\top$ versus truth $\bot$ in the first place. The bridge from
row-conditional rates to outright cell probabilities is exactly the
joint distribution we are deferring. \Cref{sec:bernoulli-bool-formal}
takes the bridge.

\subsection*{Where the chapter goes from here}

We now have the conceptual picture of the four cells but no data on how
they actually fill. \Cref{sec:sampling-rates} takes the next step:
it asks how the noisy bit is sampled in practice and how the rates of
\Cref{sec:quantifying-noise} are estimated from a finite number of
trials. That returns the chapter to concrete arithmetic and sets up the
question, picked up by chapters 2 through 4, of what an algebra of these
noisy bits looks like.
